\chapter{Zusammenfassung und Ausblick}

\section{Zusammenfassung}

Ziel dieser Arbeit war die Entwicklung und Evaluation eines Smart Mirrors, der neben klassischen Informationsanzeigen auch biometrische Daten kontaktlos erfassen und darstellen kann. Im Mittelpunkt standen dabei die kamerabasierte Pulsmessung sowie eine einfache Erkennung des Gesichtsausdrucks. Dazu wurde ein Prototyp aufgebaut, der ein Display hinter einem Spionspiegel, eine Kamera, einen Raspberry Pi 4B sowie eine containerisierte Softwareumgebung miteinander verbindet. Die Verarbeitung erfolgt in einer \ac{ros}-Noetic-Umgebung innerhalb eines Docker-Containers. Die Ergebnisse werden über eine WebUI dargestellt, die über \texttt{rosbridge} mit den entsprechenden \ac{ros}-Topics verbunden ist.

Für die Pulsmessung wurde ein bestehender \ac{evm}-basierter Ansatz in die Smart-Mirror-Architektur integriert und um eine einfache Schnittstelle zur WebUI erweitert. Die berechneten Pulswerte werden über das Topic \texttt{/heart\_rate} bereitgestellt. Ergänzend wurde eine ressourcenschonende Emotionserkennung umgesetzt, die mithilfe von OpenCV Haar-Cascades zwischen den Zuständen \glqq Neutral\grqq\ und \glqq Smiling\grqq\ unterscheidet und das Ergebnis über das Topic \texttt{/mood} veröffentlicht. Damit konnte gezeigt werden, dass Pulsmessung, Gesichtsausdruckserkennung und Darstellung in einer gemeinsamen Smart-Mirror-Anwendung technisch zusammengeführt werden können.

Die Evaluation zeigt jedoch auch die Grenzen des entwickelten Systems. Im Vergleich mit einer Smartwatch als Referenzgerät erreichte der rote Farbkanal mit einem MAE von \(6{,}00\,\mathrm{BPM}\) das beste Ergebnis. Der grüne Farbkanal lag mit \(9{,}83\,\mathrm{BPM}\) deutlich darüber. In den Beleuchtungstests bewegten sich die mittleren absoluten Abweichungen zwischen \(8{,}17\,\mathrm{BPM}\) und \(8{,}85\,\mathrm{BPM}\). Die durchschnittlichen Abweichungen lagen damit in einem grundsätzlich plausiblen Bereich, einzelne Ausreißer von bis zu \(42\,\mathrm{BPM}\) zeigen jedoch, dass die Messung noch nicht durchgehend stabil ist.

Als zentrale Einflussfaktoren wurden Beleuchtung, Farbkanal, stabile Gesichtserkennung, Bewegungen der nutzenden Person und die Performance des Raspberry Pi 4B identifiziert. Besonders die hohe Rechenlast der parallelen Bildverarbeitung führte im Betrieb zu einer deutlichen Erwärmung des Systems und damit zu einer Verschlechterung der Stabilität. Die Emotionserkennung funktionierte unter guten Bedingungen grundsätzlich, wurde jedoch nur qualitativ bewertet und erlaubt daher keine Aussage über eine genaue Erkennungsrate.

Insgesamt ist der Smart Mirror als funktionsfähiger Prototyp zu bewerten. Die Arbeit zeigt, dass ein alltäglicher Spiegel technisch zu einer Plattform für kontaktlose Vitaldatenerfassung und einfache Gesichtsausdruckserkennung erweitert werden kann. Für eine verlässliche gesundheitsbezogene oder medizinische Nutzung reicht die erreichte Genauigkeit und Stabilität jedoch noch nicht aus. Der entwickelte Aufbau eignet sich damit vor allem als Demonstration und als Grundlage für weitere Untersuchungen.

\section{Ausblick}

Für eine Weiterentwicklung des Systems sollte zunächst die Evaluation der Pulsmessung systematischer durchgeführt werden. Anstelle einer Smartwatch wäre ein medizinisch geeigneteres Referenzgerät sinnvoll, um die Messgenauigkeit besser bewerten zu können. Zusätzlich sollten größere Tests mit mehreren Testpersonen, unterschiedlichen Hauttypen, Abständen, Kopfpositionen und Beleuchtungssituationen durchgeführt werden. Dadurch ließe sich genauer untersuchen, unter welchen Bedingungen die kamerabasierte Pulsmessung zuverlässig arbeitet und wann sie an ihre Grenzen stößt.

Ein weiterer wichtiger Ansatzpunkt ist die Verbesserung der Signalqualität. Dazu gehören eine konstantere Beleuchtung, eine robustere Erkennung und Verfolgung der \ac{roi} sowie eine bessere Behandlung von Bewegungsartefakten. Auch eine Plausibilitätsprüfung der berechneten Pulswerte wäre sinnvoll, damit unrealistische Werte oder starke Ausreißer nicht ungefiltert an die WebUI weitergegeben werden. Ergänzende Glättungs- und Filterverfahren könnten dazu beitragen, kurzfristige Fehlmessungen zu reduzieren und stabilere Verläufe anzuzeigen.

Auch die Hardwareplattform sollte weiter betrachtet werden. Die Beobachtungen während der Evaluation deuten darauf hin, dass der Raspberry Pi 4B bei gleichzeitiger \ac{evm}-Pulsmessung und Emotionserkennung stark belastet wird. Eine verbesserte Kühlung könnte die thermisch bedingte Reduzierung der Taktfrequenz des Prozessors verringern. Alternativ wäre der Einsatz eines leistungsfähigeren Microcontrollers zu prüfen, insbesondere wenn zukünftig komplexere Verfahren zur Pulsmessung oder Emotionserkennung verwendet werden sollen.

Die Emotionserkennung kann ebenfalls ausgebaut werden. In der vorliegenden Umsetzung wurde bewusst nur zwischen \glqq Neutral\grqq\ und \glqq Smiling\grqq\ unterschieden. Für eine umfassendere Bewertung wäre eine quantitative Messreihe notwendig, in der Erkennungsrate, Fehlklassifikationen und Reaktionszeit dokumentiert werden. Langfristig könnte außerdem ein mehrklassiges Modell eingesetzt werden, das weitere Gesichtsausdrücke erkennt. Dabei müsste jedoch weiterhin berücksichtigt werden, dass Gesichtsausdrücke keine eindeutige Aussage über den tatsächlichen emotionalen Zustand einer Person bieten.

Da die Verarbeitung der Kamera- und Messdaten im entwickelten System lokal erfolgt, ist bereits eine wichtige Grundlage für einen datensparsamen Betrieb gegeben. Für eine zukünftige Version wäre vor allem sinnvoll, diese lokale Verarbeitung für die Nutzer transparenter zu machen. Dazu könnten beispielsweise ein klar sichtbarer Kamerastatus sowie einfache Möglichkeiten zur Deaktivierung der Kamera oder einzelner Funktionen beitragen. Für den praktischen Einsatz im häuslichen Umfeld stehen damit vor allem die weitere Verbesserung der technischen Stabilität, der Messqualität und der nachvollziehbaren Bedienung im Vordergrund.
